
% ════════════════════════════════════════════
% 模型建立（5.2.3.模型建立.tex，三级子块）
% ════════════════════════════════════════════
\subsubsection{两阶段随机规划与紧急结算模型}

问题二的核心困难在于信息的时间结构：微网在每天 0:00 制定当日购电与充放电计划时，既不知道当天的负载曲线 $L_{d,k}$，也不知道当天的光伏出力 $S_{d,k}$（如假设3所述）；而一旦实际功率不足，缺口必须以交易电价的 5 倍紧急购电。这意味着本问不能沿用问题一"已知当日全部数据再求解"的确定性框架。若在规划时直接代入附件二的真实值，由于计划购电与充放电可自由调节，最优解必然逐点覆盖真实缺口，供电缺口 $\delta_{d,k}$ 恒为零、紧急购电费用恒为零，"5 倍电价"这一惩罚机制将完全失效，模型也就失去了对预测偏差的刻画能力，本质上属于对已知数据的"事后拟合"。

因此，本节建立\textbf{"先预测、再决策、后结算"}的三段式模型：每日 0:00 仅利用截至前一日的全部历史数据，预测当日 144 个间隔的负载与光伏；由历史预测残差的经验分布构造有限的偏差情景，用两阶段随机规划在第一阶段（当日 0:00 一次性的确定性计划购电量与充放电功率）与第二阶段（情景依赖的紧急购电）之间进行权衡，得到兼顾预期费用与偏差风险的当日计划；待当日实际功率到达后，按真实缺口结算紧急购电费用。三段之间以"预测信息 $\to$ 计划决策 $\to$ 实际结算"的顺序严格单向衔接，不使用任何当日未来的真实信息，从而避免了上述确定性模型的缺陷。具体过程如下。

\textbf{步骤1：预测模型的标定}

设第 $d$ 天的历史数据为 $\{L_{\tau,k},S_{\tau,k}\}_{\tau<d}$，当日预测为
\begin{equation}
    \hat L_{d,k}=\mathrm{lev}^{L}_{d}\cdot \mathrm{shp}^{L}_{d,k},\qquad
    \hat S_{d,k}=\mathrm{lev}^{S}_{d}\cdot \mathrm{shp}^{S}_{d,k},\qquad k=0,1,\dots,143
\end{equation}
即"日水平 $\times$ 日内形状"的乘积形式。负载具有明显的星期效应，故先按"工作日/周末"将历史日分为两类，取与目标日同类的最近 $k_s=8$ 个历史日构造归一化日形状，并对其日均水平作近期加权：
\begin{equation}
    \mathrm{shp}^{L}_{d,k}=\frac{1}{k_s}\sum_{\tau\in\mathcal{K}^{L}_{d}}
    \frac{L_{\tau,k}}{\bar L_{\tau}},\qquad
    \bar L_{\tau}=\frac{1}{144}\sum_{k=0}^{143}L_{\tau,k}
\end{equation}
\begin{equation}
    \mathrm{lev}^{L}_{d}=\frac{\sum_{i=1}^{k_\ell} w_i\,\bar L_{\tau_i}}{\sum_{i=1}^{k_\ell} w_i},
    \qquad w_i=1+\frac{i-1}{k_\ell-1}
\end{equation}
其中 $\tau_1<\dots<\tau_{k_\ell}$ 为同类历史日中最近的 $k_\ell=2$ 天，权重 $w_i$ 由 1 线性递增至 2，使越近的日期权重越大。光伏的日内形态对星期不敏感，故直接取最近 $k_s=4$ 天的归一化形状均值与最近 $k_\ell=4$ 天的加权日均水平相乘得到 $\hat S_{d,k}$。

上述标定采用\textbf{扩张窗口}：第 $d$ 天的模型只使用 $d$ 之前的数据，随着日期推进，可供标定的历史样本逐日增加。附件二的完整记录自 2025 年 1 月 1 日开始，1 月的 31 天数据正是用于标定上述预测模型的历史信息，因此本问的滚动优化统计区间取 2025 年 2 月 1 日至 12 月 31 日，共 334 天；1 月作为预测模型的预热与标定期，其结果不计入全年费用汇总。这一处理也与题目所给结果模板从 2 月 1 日起算的要求相一致。

\textbf{步骤2：预测误差分布与偏差情景构造}

预测必然存在误差，而紧急购电正是对误差的兜底。为刻画误差，对每一天计算预测的相对残差：
\begin{equation}
    \varepsilon^{L}_{d,k}=\frac{L_{d,k}-\hat L_{d,k}}{\max\{\hat L_{d,k},1\}},\qquad
    \varepsilon^{S}_{d,k}=\frac{S_{d,k}-\hat S_{d,k}}{\max\{\hat S_{d,k},1\}}\quad(\hat S_{d,k}>50)
\end{equation}
其中光伏仅在预测出力大于 $50$ kW 时计入残差，以避免夜间零出力时刻的相对误差被过度放大。将滚动过程中已经发生的残差逐点存入残差池 $\mathcal{R}^{L}$、$\mathcal{R}^{S}$，并取经验分位数标定三档偏差因子：
\begin{equation}
    \lambda^{L}_{a}=1+Q_{q_a}\!\left(\mathcal{R}^{L}\right),\qquad
    \lambda^{S}_{b}=1+Q_{q_b}\!\left(\mathcal{R}^{S}\right),\qquad
    (q_1,q_2,q_3)=(0.1,0.5,0.9)
\end{equation}
其中 $Q_{q}(\cdot)$ 为经验分位数函数，$\lambda^{L}$ 与 $\lambda^{S}$ 分别裁剪至 $[0.7,1.3]$ 与 $[0.5,1.5]$ 以保证情景合理；当有效样本不足 50 个时退回到先验因子 $(0.95,1.00,1.05)$。由实测残差标定得到的三档因子均值为：负载 $\lambda^{L}=(0.9526,0.9986,1.0435)$，光伏 $\lambda^{S}=(0.9073,1.0139,1.1301)$，可见光伏预测偏差明显大于负载，这与光伏出力受云量、辐照随机性影响的物理事实相符。

将负载与光伏的三档因子两两组合，得到 9 个偏差情景 $\omega_s=(\lambda^{L}_{a},\lambda^{S}_{b})$，$a,b\in\{1,2,3\}$，情景概率取两档边缘概率的外积：
\begin{equation}
    L^{s}_{d,k}=\lambda^{L}_{a}\,\hat L_{d,k},\qquad
    S^{s}_{d,k}=\lambda^{S}_{b}\,\hat S_{d,k},\qquad
    \pi_s=w_a\,w_b,\qquad (w_1,w_2,w_3)=(0.2,0.6,0.2)
\end{equation}
即在"负载低/中/高 $\times$ 光伏低/中/高"共 9 种偏差组合下刻画当日可能出现的实际功率。

\textbf{步骤3：两阶段随机规划决策}

在第一阶段，微网于当日 0:00 一次性确定计划购电功率 $g_{d,k}$、充电功率 $c_{d,k}$、放电功率 $d_{d,k}$ 与储电量 $E_{d,k}$，这些决策\textbf{不依赖情景}，即在 0:00 只能得到与情景无关的唯一计划；在第二阶段，当日实际功率揭晓后，各情景下的供电缺口 $\delta^{s}_{d,k}$ 按 5 倍电价结算。当日决策模型为
\begin{equation}
    \min\; Z_d=\underbrace{\sum_{k=0}^{143} g_{d,k}\,p_k\,\Delta t}_{\text{计划购电费}}
    +\underbrace{\sum_{s=1}^{9}\pi_s\sum_{k=0}^{143}5\,p_k\,\delta^{s}_{d,k}\,\Delta t}_{\text{预期紧急购电费}}
\end{equation}
\begin{equation}
    \delta^{s}_{d,k}\geq L^{s}_{d,k}-\left(g_{d,k}+S^{s}_{d,k}+d_{d,k}-c_{d,k}\right),\qquad
    \delta^{s}_{d,k}\geq 0,\qquad \forall s,\ \forall k
    \label{eq:q2_deltacon}
\end{equation}
式~\eqref{eq:q2_deltacon}~中的取最大值关系 $\delta^{s}_{d,k}=\max\{\cdot,0\}$ 已由"下界约束 $+$ 目标函数中系数为正"共同线性化：$\delta^{s}_{d,k}$ 在最优点处必然取到下界，无需引入整数变量。储电量递推与容量、功率约束为
\begin{equation}
    E_{d,k+1}=E_{d,k}+\eta\,c_{d,k}\,\Delta t-\frac{1}{\eta}\,d_{d,k}\,\Delta t
    \label{eq:q2_erecur}
\end{equation}
\begin{equation}
    E_{\min}\leq E_{d,k}\leq E_{\max},\qquad
    0\leq c_{d,k},\,d_{d,k}\leq P_{\max},\qquad g_{d,k}\geq 0
\end{equation}
其中 $p_k$ 为附件一给出的与日期无关的分时电价，$\eta=0.9$ 为充放电效率，$P_{\max}=5000$ kW，$E_{\min}=1200$ kWh，$E_{\max}=10800$ kWh，$\Delta t=1/6$ h。式~\eqref{eq:q2_erecur}~对全部情景唯一，体现了储电量轨迹的物理确定性（同一储能装置只有一条实际轨迹），而 $\delta^{s}_{d,k}$ 逐情景独立，体现了第二阶段随信息揭晓而"因情制宜"的适应能力。该模型即两阶段随机线性规划的标准形式\cite{birge2011}。

\textbf{步骤4：实际运行结算与紧急购电、富余电量计算}

当日实际负载 $L_{d,k}$ 与光伏 $S_{d,k}$ 到达后，如假设6所设定，储能按第一阶段的计划充放电功率执行，具体为：当实际功率不足时，缺口部分以 5 倍电价紧急购电补足
\begin{equation}
    \delta_{d,k}=\max\left\{L_{d,k}-\left(g_{d,k}+S_{d,k}+d_{d,k}-c_{d,k}\right),\,0\right\}
    \label{eq:q2_delta}
\end{equation}
\begin{equation}
    Z^{\mathrm{em}}_{d}=\sum_{k=0}^{143}5\,p_k\,\delta_{d,k}\,\Delta t
\end{equation}
当实际功率富余（如光伏出力超过预测）时，储能按计划功率继续充电直至容量上限，超出 $E_{\max}$ 的部分无法消纳，形成弃电
\begin{equation}
    su_{d,k}=\max\left\{S_{d,k}+g_{d,k}+d_{d,k}-c_{d,k}-L_{d,k},\,0\right\}\Delta t
    \label{eq:q2_curtail}
\end{equation}
当日计划购电费为 $Z^{\mathrm{plan}}_{d}=\sum_k p_k\,g_{d,k}\,\Delta t$。式~\eqref{eq:q2_delta}~与式~\eqref{eq:q2_curtail}~表明：在三段式框架下，预测偏差既可能表现为需要 5 倍价格兜底的供电缺口，也可能表现为无法消纳的光伏富余，二者均由实际值与预测值的偏差驱动，这正是该模型能够刻画紧急购电机制的关键。

\textbf{步骤5：储电量日间传递与全年费用汇总}

储电量在全年范围内连续传递，当日末值即次日初值：
\begin{equation}
    E_{d+1,0}=E_{d,144},\qquad E_{0,0}=6000
\end{equation}
与问题一要求的"日循环（0:00 与 24:00 储电量相等）"不同，问题二不设终端约束，储电量可跨日留存。但由于电价曲线日间完全相同，跨日留存只能在次日相同电价下使用，还需额外承担约 $19\%$ 的往返效率损耗（$1-0.9^2=0.19$），因此不存在跨日套利空间，最优解仍在日内完成"低价时段充电、高价时段放电"。实测的四个指定日期 0:00 与 24:00 储电量均恰为下限 $1200.00$ kWh（见表\ref{tab:结果2-2a}～表\ref{tab:结果2-2d}），说明日内充放电空间已被充分利用，日间不存在需要额外协调的储电量耦合，逐日滚动求解在最优值上与全年统一求解等价\cite{sioshansi2010,zhoulin2011,sangbingyu2014}。

综上，问题二全年的总购电费用为
\begin{equation}
    Z=\sum_{d=1}^{334}\left(Z^{\mathrm{plan}}_{d}+Z^{\mathrm{em}}_{d}\right)
    =\sum_{d=1}^{334}\sum_{k=0}^{143}p_k\,g_{d,k}\,\Delta t
    +\sum_{d=1}^{334}\sum_{k=0}^{143}5\,p_k\,\delta_{d,k}\,\Delta t
\end{equation}
模型按天递推求解 $334$ 次，每次求解一个含 $13\times144+1=1873$ 个决策变量、$9\times144$ 条缺口约束的线性规划，全年总耗时约 $10.6$ s，既保证了计算效率，又通过储电量的逐日传递与残差池的逐日扩充实现了"预测精度随数据积累而提高"的滚动学习效果。

本节模型可概括为"先预测、再决策、后结算"的三段式闭环：扩张窗口标定负载与光伏的预测模型并给出偏差情景，两阶段随机规划据此制定当日的充放电与购电计划，实际功率到达后按 5 倍电价结算紧急购电与弃电。串联三者的两条纽带是储电量的逐日传递与残差池的逐日扩充，前者保证储能轨迹在物理上可行，后者使偏差分布随数据积累而更新。以该模型为基础，下一节的求解将给出全年逐日的购电量、购电费与紧急购电分布。
